% =============================================================================
% Anexos (iteración 3c, 7-nov-2026). Van después de las referencias (guía del curso MMC102).
% =============================================================================

% -----------------------------------------------------------------------------
\chapter{Demostraciones técnicas}
\label{ap:demostraciones}

Este anexo reúne las demostraciones del \cref{cap:marco} que, por su extensión, cortarían la
lectura del capítulo. Todas son propias. Los enunciados están en el capítulo y no se repiten.

\section{Ganancia de la primera diferencia}
\label{ap:dem-ganancia}

Se demuestra la \cref{prop:ganancia}, que sustenta la decisión de analizar las series sin
transformar (\cref{sec:epistemologia}).

\begin{proof}[Demostración de la \cref{prop:ganancia}]
Por la identidad
\[
  \cos\xi_1-\cos\xi_2=-2\sin\frac{\xi_1+\xi_2}{2}\sin\frac{\xi_1-\xi_2}{2},
\]
con $\xi_1=\vartheta t+\psi_0$ y $\xi_2=\vartheta(t-1)+\psi_0$,
\[
  \Delta x_t=-2\varrho\sin\frac{\vartheta}{2}\,\sin\Bigl(\vartheta t+\psi_0-\frac{\vartheta}{2}\Bigr)
  =2\varrho\sin\frac{\vartheta}{2}\,\cos\Bigl(\vartheta t+\psi_0-\frac{\vartheta}{2}+\frac{\pi}{2}\Bigr),
\]
porque $-\sin\xi=\cos(\xi+\pi/2)$ para todo $\xi$. Como $P\ge2$, $\vartheta/2=\pi/P\in(0,\pi/2]$ y
$\sin(\pi/P)>0$, de modo que la amplitud es $2\varrho\sin(\pi/P)=\gamma(P)\varrho$. Además,
$\abs{1-e^{-\mathrm{i}\vartheta}}^2=(1-\cos\vartheta)^2+\sin^2\vartheta=2-2\cos\vartheta
=4\sin^2(\vartheta/2)$, lo que prueba la segunda igualdad de \cref{eq:ganancia}.
\textup{(a)} Por \cref{eq:ganancia}, $x$ y $\Delta x$ son de la forma
$y_t=\varrho_y\cos(\vartheta t+\psi_y)$, con $(\varrho_y,\psi_y)=(\varrho,\psi_0)$ y
$(\varrho_y,\psi_y)=\bigl(\gamma(P)\varrho,\,\psi_0+(\pi-\vartheta)/2\bigr)$, respectivamente.
Para $k\in\{1, 2\}$ y $\xi\in\RR$, si $e^{\mathrm{i}k\vartheta}\ne1$, la suma geométrica da
\[
  \Bigl\lvert\sum_{t=t_0+1}^{t_0+T}e^{\mathrm{i}(k\vartheta t+\xi)}\Bigr\rvert
  =\frac{\abs{e^{\mathrm{i}k\vartheta T}-1}}{\abs{e^{\mathrm{i}k\vartheta}-1}}
  \le\frac{2}{\abs{e^{\mathrm{i}k\vartheta}-1}},
\]
y la suma es nula si además $k\vartheta T\in2\pi\mathbb{Z}$. Tomando la parte real y usando
$y_t^2=\frac{\varrho_y^2}{2}+\frac{\varrho_y^2}{2}\cos(2\vartheta t+2\psi_y)$ se distinguen dos
casos. \emph{Caso $P>2$.} Entonces $\vartheta\in(0,\pi)$ y $2\vartheta\in(0,2\pi)$, de modo que
$e^{\mathrm{i}\vartheta}\ne1$ y $e^{2\mathrm{i}\vartheta}\ne1$. Si $P$ es entero, $\vartheta P=2\pi$ y
$2\vartheta P=4\pi$, las dos sumas sobre $T=P$ meses son nulas y
$\prom{y}_{t_0,P}=0$, $\prom{y^2}_{t_0,P}=\varrho_y^2/2$, luego
$\vart_{t_0,P}(y)=\varrho_y^2/2$ exactamente. Para $P$ real cualquiera, las sumas están acotadas
uniformemente en $t_0$ y $T$, de modo que $\prom{y}_{t_0,T}\to0$ y
$\prom{y^2}_{t_0,T}\to\varrho_y^2/2$, luego $\vart_{t_0,T}(y)\to\varrho_y^2/2$ sin depender de
$t_0$. En ambos sentidos la varianza pasa de $\varrho^2/2$ a $\gamma(P)^2\varrho^2/2$.
\emph{Caso $P=2$.} Entonces $\vartheta=\pi$, $e^{2\mathrm{i}\vartheta}=1$ y el promedio de
$\cos^2$ ya no es $1/2$: $x_t=\varrho(-1)^t\cos\psi_0$, de modo que $x_t^2=\varrho^2\cos^2\psi_0$
para todo $t$ y $\prom{x}_{t_0,T}$ vale $0$ si $T$ es par y tiende a $0$ en general; luego
$\vart_{t_0,2}(x)=\varrho^2\cos^2\psi_0$ y $\vart_{t_0,T}(x)\to\varrho^2\cos^2\psi_0$. Además
$x_{t-1}=-x_t$, de donde $\Delta x_t=2x_t=\gamma(2)\,x_t$ y
$\vart_{t_0,T}(\Delta x)=\gamma(2)^2\vart_{t_0,T}(x)$ para todo $t_0$ y todo $T$.
\textup{(b)} $\pi/P$ decrece en $P$ dentro de
$(0,\pi/2]$, donde el seno es creciente; y $\sin x\sim x$ cuando $x\to0$. \textup{(c)}
$2\sin(\pi/P)<1$ equivale a $\sin(\pi/P)<1/2$, es decir, a $\pi/P<\pi/6$ dentro de
$(0,\pi/2]$, o sea, $P>6$.
\end{proof}

\section{Tamaño efectivo de los pesos de la ventana}
\label{ap:dem-kish}

Se demuestra la \cref{prop:kish}, que justifica usar el tamaño efectivo, y no el ancho, como
escala de los umbrales por ventana (\cref{def:umbral-ventana}).

\begin{proof}[Demostración de la \cref{prop:kish}]
(a) Por la desigualdad de Cauchy-Schwarz,
$1=\bigl(\sum_u\omega_u\cdot1\bigr)^2\le\nw\sum_u\omega_u^2$, es decir, $\nomega\le\nw$, con
igualdad si y solo si $(\omega_u)$ es proporcional a $(1,\dots,1)$, o sea, si $\omega_u=1/\nw$.
Como $0\le\omega_u\le1$, se tiene $\omega_u^2\le\omega_u$, con igualdad si y solo si
$\omega_u\in\{0, 1\}$; sumando, $\sum_u\omega_u^2\le1$, es decir, $\nomega\ge1$, con igualdad si y
solo si todos los pesos valen cero o uno, lo que, con suma uno, deja un único peso igual a uno.
(b) Con $\nu=0$ los pesos son uniformes y se aplica (a). Con $\nu>0$, $\tilde\omega_u$ es
estrictamente decreciente en $\abs{u-(\nw+1)/2}$. Si $\nw\ge3$, esa distancia vale $(\nw-1)/2$ en
$u=1$ y $(\nw-3)/2$ en $u=2$, de modo que $\tilde\omega_2>\tilde\omega_1$; los pesos no son todos
iguales y (a) da la desigualdad estricta. (c) La \cref{eq:pesos} solo contiene $u$, $\nw$ y $\nu$.
(d) Sea $\xi_u=a_\nu(2u-\nw-1)/(\nw-1)$, de modo que $\tilde\omega_u=e^{-\xi_u^2/2}$ y los $\xi_u$ son
$\nw$ puntos equiespaciados de $-a_\nu$ a $a_\nu$, con separación $2a_\nu/(\nw-1)$. Para $k\in\{1, 2\}$ sea
$S_k=\frac{2a_\nu}{\nw-1}\sum_u e^{-k\xi_u^2/2}$. Son sumas de Riemann de funciones continuas en
$[-a_\nu,a_\nu]$, salvo un término de orden $1/\nw$, de modo que $S_1\to I_1=\int_{-a_\nu}^{a_\nu}e^{-x^2/2}dx$ y
$S_2\to I_2=\int_{-a_\nu}^{a_\nu}e^{-x^2}dx$. Como $\nomega=(\sum_u\tilde\omega_u)^2/\sum_u\tilde\omega_u^2$,
\[
  \frac{\nomega}{\nw}=\frac{\nw-1}{\nw}\cdot\frac{S_1^2}{2a_\nu\,S_2}\;\longrightarrow\;
  \frac{I_1^2}{2a_\nu\,I_2}=c(\nu).
\]
Por la desigualdad de Cauchy-Schwarz en $[-a_\nu,a_\nu]$,
$I_1^2=\bigl(\int e^{-x^2/2}\cdot1\,dx\bigr)^2\le2a_\nu\int e^{-x^2}dx=2a_\nu\,I_2$, con igualdad solo si
$e^{-x^2/2}$ fuera constante en $[-a_\nu,a_\nu]$, lo que no ocurre con $a_\nu>0$. Luego $c(\nu)<1$, y
$c(\nu)>0$ porque $I_1>0$.
\end{proof}

\section{No negatividad del estimador HAC}
\label{ap:dem-hac}

Se demuestra el \cref{lem:hac}, que garantiza que el estadístico de la deriva de la
\cref{def:hac} esté bien definido.

\begin{proof}[Demostración del \cref{lem:hac}]
Al desarrollar el cuadrado, el producto $\dot x_s\dot x_{s'}$ con $\abs{s-s'}=k\le\bar m-1$
aparece una vez por cada $t$ con $s,s'\in\{t-\bar m+1,\dots,t\}$, es decir, $\bar m-k$ veces. Por
tanto, la suma vale
$\bar m\sum_s\dot x_s^2+2\sum_{k=1}^{\bar m-1}(\bar m-k)\sum_s\dot x_s\dot x_{s-k}
=n'\bar m\bigl[\hat\Gamma_0+2\sum_{k=1}^{\bar m-1}(1-k/\bar m)\hat\Gamma_k\bigr]$, que es
$n'\bar m\,\hat\sigma^2_{\mathrm{HAC},j}$. Si la suma es nula, todas las sumas móviles lo son; la
primera, $t=1$, da $\dot x_1=0$, y por inducción $\dot x_t=0$ para todo $t$.
\end{proof}

\section{Validez de la versión simétrica del criterio de OE4}
\label{ap:dem-oe4}

Se demuestra la \cref{prop:oe4}, que da el nivel del criterio de ancho de ventana bajo
intercambiabilidad (\cref{sup:intercambiable}).

\begin{proof}[Demostración de la \cref{prop:oe4}]
Cada $\tilde\varpi_j$ es la misma función de $\mathcal{P}_j$ y del conjunto de los demás perfiles, de modo que permutar
los perfiles permuta los $\tilde\varpi_j$; por el supuesto, $(\tilde\varpi_0,\dots,\tilde\varpi_B)$ es intercambiable. Para
$k\in\{1,\dots,B+1\}$ sea $\mathcal{J}=\{j:\#\{c\ne j:\tilde\varpi_c\ge\tilde\varpi_j\}\le k-1\}$. Si $\mathcal{J}$ tuviera $k+1$ elementos
o más, el de menor $\tilde\varpi$ en $\mathcal{J}$ tendría al menos $k$ índices distintos con $\tilde\varpi$ mayor o igual, contra
la definición de $\mathcal{J}$; luego $\#\mathcal{J}\le k$. Por la intercambiabilidad, cada índice pertenece a $\mathcal{J}$ con
la misma probabilidad, de modo que $\Pr(0\in\mathcal{J})=\mathbb{E}[\#\mathcal{J}]/(B+1)\le k/(B+1)$. Como $0\in\mathcal{J}$
equivale a $\widetilde{\pval}\le k/(B+1)$, y $\widetilde{\pval}$ solo toma valores $k/(B+1)$, se sigue
$\Pr(\widetilde{\pval}\le\alpha_0)\le\lfloor\alpha_0(B+1)\rfloor/(B+1)\le\alpha_0$.
\end{proof}

% -----------------------------------------------------------------------------
% =============================================================================
% Anexo B · Robustez con las series transformadas (decisión D15)
% Fuentes: data/results/final/anexo/ (05c_*.csv y 09b_discriminacion_episodios.csv),
% notebooks 05c y 09b, R/ciclo_financiero_final.Rmd (anexo R). Cifras solo por macro.
% =============================================================================
\chapter{Robustez con las series transformadas}
\label{ap:robustez}

Por la decisión D15, el análisis principal usa las series sin transformar y las transformadas
pasan a este anexo con dos fines: sustentar la decisión de no diferenciar
(\cref{sec:ap-no-diferenciar}) y evaluar la detección de episodios fuera de muestra, donde solo la
vía transformada es informativa (\cref{sec:ap-deteccion}). Las series transformadas cubren
\nMesesTransformadas{} meses, de \fechaInicioTransformadas{} a \fechaFinDatos{}, con las mismas
\nSeries{} series del análisis principal. Este anexo no reabre la Parte~C de la decisión D2, que
sigue pendiente de discutir con el asesor (\cref{sec:pend-parte-c}).

\section{Sustento de no diferenciar}
\label{sec:ap-no-diferenciar}

La transformación de las series combina dos pasos: el logaritmo, en las series de magnitudes
positivas, y la primera diferencia. El notebook \texttt{05c} los separa y repite el PCA después
de cada uno.

\paragraph{El logaritmo no cambia la estructura.} Tras aplicar solo el logaritmo, la congruencia
de Tucker (\cref{def:tucker}) de los tres primeros ejes con los de las series sin transformar es
\congruenciaLogPCUno{},
\congruenciaLogPCDos{}
y \congruenciaLogPCTres{},
y la inercia del primer eje casi no cambia. La discusión sobre la transformación es, por tanto,
una discusión sobre diferenciar.

\paragraph{La diferencia sí la cambia.} Con solo la primera diferencia, la inercia del primer
eje cae de \inerciaPCUnoNiveles\,\% a
\inerciaPCUnoDiferencia{}\,\%,
y su correlación con el tiempo desaparece. La \cref{prop:ganancia} explica por qué: la primera
diferencia multiplica la varianza de una oscilación de período $P$ meses por
$\gamma(P)^2=4\sin^2(\pi/P)$, que es menor que uno para todo $P>6$ y tiende a cero como
$(2\pi/P)^2$. Una oscilación de ocho años ($P=96$) conserva solo la fracción $4\sin^2(\pi/96)$
de su varianza; el archivo \texttt{05c\_ganancia.csv} tabula la ganancia para los períodos de
interés. En los datos, la proporción de la varianza de cada serie en oscilaciones de dos años o
más pasa, en mediana, de
\espectroLentoMedianaNivelesSinTendencia\,\%
en las series sin transformar y sin tendencia lineal a
\espectroLentoMedianaTransformadas\,\%
en las transformadas.

\paragraph{Lectura.} Diferenciar elimina justamente la escala de tiempo que define el ciclo
financiero (\cref{sec:ea-ciclo}). Analizar las series sin transformar conserva esa escala, a
cambio de mezclar la tendencia de profundización con la estructura de régimen, mezcla que la
regla D5 y el nulo de tendencias mixto permiten leer (\cref{sec:res-tendencia}).

\section{Detección de episodios fuera de muestra}
\label{sec:ap-deteccion}

La razón de error \emph{walk-forward} del PCA (\cref{def:walkforward}) se calcula con $d=3$ en
las dos vías. Con las series transformadas, la proporción de falsas alarmas fuera de los
episodios ampliados en 3 meses es
\falsasAlarmasTransformadasTres{}\,\%,
y la sensibilidad en la pandemia de COVID-19 es
\sensibilidadCovidTransformadasTres{}\,\%.
Con las series sin transformar, en cambio, las falsas alarmas suben a
\falsasAlarmasNivelesTres{}\,\%:
el PCA ajustado con los meses anteriores no reconstruye bien los meses nuevos porque el nivel de
las series sigue creciendo, de modo que casi todo mes parece atípico. Por eso la decisión D15
conserva la detección de episodios en la vía transformada.

Dos limitaciones acotan esta lectura. Los episodios son pocos y algunos puntuales, de modo que la
sensibilidad se calcula sobre muy pocos meses y no tiene poder estadístico. Y los episodios que
la vía transformada detecta son choques abruptos, como la pandemia, no cambios lentos de
estructura, que son los que interesan a las hipótesis HT1 y HT3. La complementariedad entre ambas
vías es precisamente la propuesta de la Parte~C de D2, que se discutirá con el asesor.


% -----------------------------------------------------------------------------
\chapter{Notación}
\label{ap:notacion}

La tabla siguiente reúne los símbolos de la tesis, agrupados por método. La notación es única para
todo el documento y se define en el archivo \texttt{tesis\_latex/notacion.tex}; antes de
introducir un símbolo se verifica que no exista. Las abreviaturas se listan en las páginas
preliminares.

\listadesimbolos

% -----------------------------------------------------------------------------
% =============================================================================
% Anexo D · Reproducibilidad: datos congelados (D4), paridad Python–R (D6) y orden de ejecución
% Fuentes: decisiones.qmd (D4, D6), CLAUDE.md (comandos), data/results/final/R_paridad_python.csv.
% Las huellas md5 son identificadores de D4 (clase «Identificadores» de la lista de excepciones).
% =============================================================================
\chapter{Reproducibilidad}
\label{ap:reproducibilidad}

Este anexo documenta cómo se reproduce el análisis. Complementa la
\cref{sec:met-implementacion}, que describe la implementación y la paridad.

\section{Datos congelados}
\label{sec:ap-datos}

Por la decisión D4, la tesis usa la versión de los datos de \fechaInicioNiveles{} a
\fechaFinDatos{} versionada en el repositorio el 16 de octubre de 2026. El código comprueba la
huella md5 de cada archivo antes de leerlo y se detiene si no coincide
(\cref{tab:md5}). Cualquier actualización de los datos requiere una nueva decisión, porque el BCRP
revisa el ajuste estacional de algunas series y sus valores cambian con la fecha de descarga.

\begin{table}[htbp]
\centering
\small
\caption{Archivos de datos congelados y su huella md5 (decisión D4)}
\label{tab:md5}
\begin{tabular}{@{}>{\raggedright\arraybackslash}p{0.5\textwidth}l@{}}
\toprule
Archivo (en \path{data/processed/}) & md5 \\
\midrule
\path{dataset_niveles_mensual.csv} & \texttt{8088c28cf277419388d9562722faa1ad} \\
\path{dataset_transformado_mensual.csv} & \texttt{8c49f3e55608f2651be3c1dee4236194} \\
\path{dataset_niveles_estandarizado_mensual.csv} & \texttt{153638578212ecc820ec9223087c080a} \\
\path{dataset_estandarizado_mensual.csv} & \texttt{43a01dead7ffb96e8b3b7e4d2799f15d} \\
\path{series_bcrp_consolidado.csv} & \texttt{61e6a6bc92eee23c6c476f096d402edb} \\
\bottomrule
\end{tabular}

\smallskip
\begin{minipage}{0.95\textwidth}\footnotesize
Fuente: \texttt{decisiones.qmd}, decisión D4. Las exclusiones de series (crédito total del sistema
financiero e IPC subyacente por D1; liquidez M2 por D16) se aplican al leer los archivos, que no
se modifican.
\end{minipage}
\end{table}

\section{Paridad entre Python y R}
\label{sec:ap-paridad}

Por la decisión D6, todo el análisis se programa en Python y en R, y la versión en R es la
referencia. La tabla de paridad del análisis principal
(\path{data/results/final/R_paridad_python.csv}) compara cada resultado en las dos lenguas con
la tolerancia numérica declarada en cada fila; en la versión vigente, coinciden
\filasParidadCoinciden{}
de \filasParidad{} comparaciones. El anexo de
robustez tiene su propia tabla (\path{data/results/final/anexo/R_paridad_python.csv}). Los
sorteos aleatorios los hace R con semilla fija y los escribe en archivos que Python lee, de modo
que la paridad no depende de los generadores de números aleatorios.

\section{Orden de ejecución}
\label{sec:ap-orden}

Desde la raíz del repositorio, la cadena vigente se reproduce en este orden:

\begin{enumerate}[leftmargin=*]
  \item notebooks de Python \path{05_pca}, \path{05c_anexo_no_diferenciar},
    \path{06_epca}, \path{06c_aede_grilla}, \path{07b_nulo_estructura},
    \path{07c_hfc_clases}, \path{09a_datacion_pca} y \path{09b_walkforward_pca}, con
    \texttt{jupyter nbconvert -{}-execute};
  \item el documento final en R, \path{R/ciclo_financiero_final.Rmd}, con
    \texttt{rmarkdown::render}, que escribe los archivos \path{R_*.csv} y la tabla de paridad;
  \item \texttt{python3} \path{herramientas/generar_cifras.py}, que regenera
    \path{tesis_latex/cifras.tex}; la opción \texttt{-{}-comprobar} verifica que esté al día;
  \item \texttt{latexmk -lualatex -cd} \path{tesis_latex/tesis.tex}, que compila este documento.
\end{enumerate}

Las pruebas automáticas del repositorio (\texttt{pytest} \path{tests/}) verifican las funciones de
\path{src/}, incluida la paridad del port de \texttt{aede} y \texttt{clahfac}.

\section{Estados anteriores recuperables}
\label{sec:ap-estados}

Los estados del análisis anteriores a las decisiones vigentes se conservan en el control de
versiones: el tag \texttt{v0.1-transformadas} (series transformadas como base, anterior a D1), el
tag \texttt{v0.2-niveles} (anterior a D2) y la rama \path{archivo/v0.3-trimestral} (frecuencia
trimestral, anterior a D3). El estado con la liquidez M2 y en dos vías, anterior a D15 y D16, está en el
commit \texttt{2be3665}; el tag \texttt{v0.4-17series-dos-vias} que debe marcarlo aún no se ha
creado (\cref{sec:pend-operativos}).


% -----------------------------------------------------------------------------
% =============================================================================
% Anexo E · Documentación complementaria (guía del curso MMC102)
% No se inventa nada: la relación de artículos y de eventos la completa el tesista.
% =============================================================================
\chapter{Documentación complementaria}
\label{ap:complementaria}

La guía del curso pide, como documentación complementaria, la relación de artículos publicados
o por publicar y las constancias de participación en eventos. Esta sección la completa el tesista
con la información y los documentos que correspondan; el redactor no consigna ningún dato que no
conste en el repositorio.

\section{Relación de artículos publicados o por publicar}
\label{sec:ap-articulos}

\begin{center}
\small
\begin{tabular}{@{}p{0.06\textwidth}p{0.42\textwidth}p{0.22\textwidth}p{0.2\textwidth}@{}}
\toprule
N.\textsuperscript{o} & Título y autores & Revista o evento & Estado \\
\midrule
 & \pendiente{relación de artículos: la completa el tesista} & & \\
\bottomrule
\end{tabular}
\end{center}

\section{Constancias de eventos}
\label{sec:ap-eventos}

\begin{center}
\small
\begin{tabular}{@{}p{0.06\textwidth}p{0.42\textwidth}p{0.22\textwidth}p{0.2\textwidth}@{}}
\toprule
N.\textsuperscript{o} & Evento & Fecha y lugar & Tipo de participación \\
\midrule
 & \pendiente{constancias de eventos: las completa el tesista} & & \\
\bottomrule
\end{tabular}
\end{center}

